\documentclass[11pt]{article}

\usepackage[a4paper,margin=30mm]{geometry}
\usepackage{amsmath,amssymb,amsthm,mathtools}
\usepackage[T1]{fontenc}
\usepackage{lmodern}
\usepackage{microtype}
\usepackage[hidelinks]{hyperref}\hypersetup{pdftitle={A Parity Classification for Maximum-Likelihood Thresholds of Distinctly Edge-Colored Cycles},pdfauthor={Xiaojian Zeng},pdfsubject={Maximum-likelihood thresholds of colored Gaussian graphical models}}

\newtheorem{theorem}{Theorem}
\newtheorem{corollary}[theorem]{Corollary}
\newcommand{\R}{\mathbb{R}}
\newcommand{\PSD}{\operatorname{PSD}}
\newcommand{\mlt}{\operatorname{mlt}}
\newcommand{\gcr}{\operatorname{gcr}}

\title{A Parity Classification for Maximum-Likelihood Thresholds\\
of Distinctly Edge-Colored Cycles}
\author{Xiaojian Zeng\\
\small PBC School of Finance, Tsinghua University, 100190\\
\small \texttt{xiaojianzeng@foxmail.com}}
\date{}

\begin{document}
\maketitle

\begin{abstract}
Consider the colored Gaussian graphical model on a cycle in which all
vertices have one color and every edge has its own color.  We prove that its
maximum-likelihood threshold is $2$ for every odd cycle length.  Combined
with the known even case, this gives threshold $1$ for even lengths and $2$
for odd lengths.  In particular, the five-cycle has threshold $2$, not $1$.
The proof constructs, on a nonempty open set of one-observation vectors, a
positive-semidefinite concentration matrix having that observation as its
one-dimensional kernel.  The construction is an exact diagonal congruence
to a positive-weight cycle Laplacian.
\end{abstract}

\section{Context and statement}

Colored Gaussian graphical models encode equality constraints among entries
of a concentration matrix; the edge- and vertex-symmetric models were
introduced by H\o jsgaard and Lauritzen~\cite{HojsgaardLauritzen2008}.  The
maximum-likelihood threshold for an uncolored graph was developed in
\cite{GrossSullivant2018}.  We use the colored-model formulation and kernel
criterion of Homs, Kuznetsova, and Stolz~\cite{HomsKuznetsovaStolz2026}.

Let $G_m$ be the undirected cycle on $\mathbb Z/m\mathbb Z$, where $m\geq3$.
All vertices form one color class, and the $m$ edges form $m$ separate color
classes.  Thus its concentration space consists of matrices $K(a,b)$ with
indices taken modulo $m$ and
\[
 K(a,b)_{ii}=a,\qquad
 K(a,b)_{i,i+1}=K(a,b)_{i+1,i}=b_i,
\]
all other entries being zero.  Separate edge colors mean that the parameters
$b_i$ are independent; they are not required to have distinct numerical
values.

For this family, Proposition~3.26 of
\cite{HomsKuznetsovaStolz2026} proves $\gcr(G_m)=2$ for every $m$ and
$\mlt(G_m)=1$ for even $m$.  Conjecture~3.27 there proposes threshold $1$ in
the odd case, although its Appendix Table~2 records threshold $2$ for the
three-vertex member.  The following theorem gives the opposite value for
every odd length and, in particular, settles the five-cycle case.

\begin{theorem}\label{thm:main}
For every $m\geq3$,
\[
 \mlt(G_m)=
 \begin{cases}
 1,&m\text{ even},\\
 2,&m\text{ odd}.
 \end{cases}
\]
\end{theorem}

The even assertion is the cited proposition.  We prove the odd assertion
below.  We use two further results from \cite{HomsKuznetsovaStolz2026}:
Theorem~3.4(2), which characterizes threshold one through the absence, for
almost every observation $x$, of a nonzero $K\in L_{G_m}\cap\PSD(m)$ with
$Kx=0$; and Corollary~3.7, which gives $\mlt(G_m)\leq\gcr(G_m)$.

\section{An open positive-measure obstruction}

\begin{proof}[Proof of Theorem~\ref{thm:main} for odd $m$]
Fix an odd $m\geq3$ and set the common diagonal parameter to $a=2$.  For an
observation $x=(x_0,\ldots,x_{m-1})^{\mathsf T}$, the equations
$K(2,b)x=0$ are
\begin{equation}\label{eq:kernel}
 x_{i-1}b_{i-1}+x_{i+1}b_i=-2x_i
 \qquad (i\in\mathbb Z/m\mathbb Z).
\end{equation}
Let $B_m(x)$ be the coefficient matrix of this system in
$b=(b_0,\ldots,b_{m-1})^{\mathsf T}$.  Row $i$ is supported in columns
$i-1$ and $i$, where its entries are $x_{i-1}$ and $x_{i+1}$, respectively.

There are exactly two supported permutations in the Leibniz expansion of
$\det B_m(x)$.  Indeed, if any row uses column $i-1$, bijectivity propagates
that choice around the cycle; otherwise every row uses column $i$.  The two
products are both $\prod_i x_i$, while their relative sign is the sign
$(-1)^{m-1}$ of an $m$-cycle.  Consequently
\begin{equation}\label{eq:determinant}
 \det B_m(x)=\bigl(1+(-1)^{m-1}\bigr)\prod_{i=0}^{m-1}x_i
             =2\prod_{i=0}^{m-1}x_i.
\end{equation}

Equation~\eqref{eq:kernel} therefore has the unique solution
\[
 b(x)=-2B_m(x)^{-1}x
\]
whenever every coordinate of $x$ is nonzero.  This solution is a continuous
rational function on the positive orthant.  At
$\mathbf 1=(1,\ldots,1)^{\mathsf T}$, uniqueness and
\eqref{eq:kernel} give $b_i(\mathbf1)=-1$ for all $i$.  Hence
\begin{equation}\label{eq:open-set}
 U_m=\{x\in(0,\infty)^m:b_i(x)<0\text{ for every }i\}
\end{equation}
is open and contains $\mathbf1$.  In particular, $U_m$ has positive
$m$-dimensional Lebesgue measure.

For $x\in U_m$, put
\[
 w_i=-b_i(x)x_ix_{i+1}>0,
 \qquad D=\operatorname{diag}(x_0,\ldots,x_{m-1}).
\]
Multiplication of the $i$th equation in \eqref{eq:kernel} by $x_i$ gives
\begin{equation}\label{eq:degree}
 2x_i^2=w_{i-1}+w_i.
\end{equation}
The diagonal identities \eqref{eq:degree} and the off-diagonal entries now
give the exact congruence
\begin{equation}\label{eq:congruence}
 DK(2,b(x))D=L_w,
\end{equation}
where $L_w$ is the Laplacian of the cycle with positive edge weights $w_i$.
Equivalently, for every $y\in\R^m$,
\begin{equation}\label{eq:squares}
 y^{\mathsf T}K(2,b(x))y
 =\sum_{i=0}^{m-1}w_i
 \left(\frac{y_i}{x_i}-\frac{y_{i+1}}{x_{i+1}}\right)^2.
\end{equation}
Because the cycle is connected, \eqref{eq:squares} proves
\[
 K(2,b(x))\succeq0,\qquad
 \ker K(2,b(x))=\R x,\qquad
 \operatorname{rank}K(2,b(x))=m-1.
\]

Thus every $x$ in the positive-measure set $U_m$ is annihilated by a
nonzero matrix in $L_{G_m}\cap\PSD(m)$.  The kernel characterization cited
above rules out $\mlt(G_m)=1$.  On the other hand,
$\mlt(G_m)\leq\gcr(G_m)=2$, so $\mlt(G_m)=2$.
\end{proof}

\begin{corollary}
For the five-cycle with one vertex color and pairwise distinct edge colors,
$\mlt(G_5)=2$.  Hence the assertion $\mlt(G_5)=1$ is false on a
positive-measure, rather than merely exceptional, set of observations.
\end{corollary}

\section{Scope, reproducibility, and assistance}

The argument is pointwise for each fixed odd $m$; it does not assert a
neighborhood in \eqref{eq:open-set} that is uniform in $m$.  For even $m$,
the determinant in \eqref{eq:determinant} vanishes identically, and the
present construction is not a proof of the even case.  The parity
classification uses the previously proved even result exactly as stated
above.  No claim of bibliographic priority beyond this comparison is made.

A supplementary Python/SymPy verifier reconstructs $B_m(x)$ directly from
\eqref{eq:kernel} and checks the determinant support, parity, exact rational
kernel equations, Laplacian congruences, and rank certificates in finite
test dimensions.  These computations audit displayed exact identities; the
arbitrary-$m$ conclusion rests on the proof above.  This use of exact
arithmetic as an audit rather than as a replacement for proof is also found
in finite positive-semidefinite matrix certification, for example
\cite{ZengLiuRatnavelu2026}.

An AI coding assistant was used to organize the proof, prepare and run the
symbolic checks, screen cited records, and draft the manuscript.  Those uses
do not constitute external peer review; all mathematical dependencies and
the full dimension-uniform argument are stated in the article.

\bibliographystyle{plain}
\bibliography{references}

\end{document}
